\documentclass{ximera}
\title{Limits and Derivatives, Part 1}


\newcommand{\pskip}{\vskip 0.1 in}

\begin{document}
\begin{abstract}
Limits in context.
\end{abstract}
\maketitle


\pskip

\section{Limits and Speed}

\begin{example}  \label{Ex:9u4dfg9}
A rock is dropped from rest near the surface of the moon. The function 
\[
      s =  f(t) = 0.8 t^2 \, , \, 0 \leq t \leq 10,
\]
expresses the distance (in meters) the rock falls in terms of the number of seconds since it was dropped.

\begin{enumerate}
\item Find a function $v=g(t)$ that gives the average speed of the rock (in meters/sec) between times $t=3$ and $t=b$ seconds. State the domain of $g$.

\item Interpret the average speed as a slope in the worksheet below. 

\begin{onlineOnly}
    \begin{center}
\desmos{tcti9t3iub}{900}{600}
\end{center}
\end{onlineOnly}

Desmos activity available at
\href{https://www.desmos.com/calculator/tcti9t3iub}{151: Moon Rock 1}

\item Drag point $B$ close to point $A$ and approximate the rock's speed $3$ seconds after being dropped.

\item Use the algebra of limits to find the exact speed of the rock three seconds after being dropped.

\item Find an expression for the rock's average speed between times $t=a$ and $t=b$.

\item Use the algebra of limits to find the speed of the rock at time $t=a$ seconds after being dropped. Compare this with the rock's average speed during the first $a$ seconds of its fall. Interpet these speeds geometrically on the graph below.

\begin{onlineOnly}
    \begin{center}
\desmos{v42lwmrlag}{900}{600}
\end{center}
\end{onlineOnly}

Desmos activity available at
\href{https://www.desmos.com/calculator/v42lwmrlag}{151: Moon Rock 2}


\end{enumerate} 

\end{example}

\section{Limits and Gas Mileage}
\begin{example} \label{Ex:9sd8gfs}
The function 
\[
   G = f(s) = \frac{2}{5} +\frac{1}{5000}(40(s+2)^2 - (s+2)^3) , \, 0\leq s \leq 23 ,
\]
expresses the number of gallons of gas in your car in terms of your distance from home. The distance is measured in miles along your route. 

\begin{onlineOnly}
    \begin{center}
\desmos{xzkknfpkw3}{900}{600}
\end{center}
\end{onlineOnly}

Desmos activity available at
\href{https://www.desmos.com/calculator/xzkknfpkw3}{151: Gas as a Function of Distance}

\begin{enumerate}
\item Use the graph of the function $f$ shown above to determine if you are driving toward or away from home. Explain your reasoning.

\item Find your average gas mileage (in miles/gallon) between distances $s=8$ and $s=18$ miles from home.

\item Use the graph to approximate your gas mileage at the moment you are $8$ miles from home. Do this by zooming in on the appropriate point.

\item Find a simplified expression for your average gas mileage (in miles/gal) between distances $s=8$ and $s=b$ miles from home. 

\emph{Hint:} Do \emph{not} simplify the expression 
\[
     f(8) = \frac{2}{5} +\frac{1}{5000}(40(10)^2 - (10)^3)
\]
for $f(8)$. Instead, leave as is and factor the difference of two squares 
\[
     A^2 - B^2 = (A+B)(A-B)
\]
and the difference of two cubes
\[
       A^3 - B^3 = (A-B)(A^2 + AB + B^2).
\]

\item Use your expression from part (d) and the algebra of limits to determine your exact gas mileage at the moment you are $8$ miles from home.

\item Find a simplified expression for your average gas mileage (in miles/gal) between distances $s=a$ and $s=b$ miles from home.  Use this expression and the algebra of limits to find an expression for your gas mileage at the moment your are $a$ miles from home.

\end{enumerate}



\end{example}




\end{document}